\section{A partition theorem for stretched data}\label{sec:interface}

The passage from scalar curvature to a map of controlled fiber diameter
uses the partition construction of \cite{OpenAIWidth}.  We state the
precise consequence needed here.  Its distances are measured in two
metrics: a metric $g$ fixes the scale of the eventual fibers, while a
larger metric $h$ separates prescribed pairs of subsets.  The work of
the subsequent sections is to construct such an $h$ from the spectral
hypothesis.  Here $h\ge g$ means quadratic-form domination:
$h_x(v,v)\ge g_x(v,v)$ for every $x\in M$ and $v\in T_xM$.

For a smooth metric $h$ and a finite family of smooth real functions
$p_\alpha$, the scalar curvature of the diagonal circle metric
\[
 h+\sum_\alpha e^{2p_\alpha}\,\dd\theta_\alpha^2
\]
is given by the standard multiple-warp formula
\cite[Section~1.A]{Gromov2024}:
\begin{equation}\label{interface:scalar}
 R(h,(p_\alpha))
 =\Scal_h-2\Delta_h P-|\dd P|_h^2
                  -\sum_\alpha|\dd p_\alpha|_h^2,
 \qquad P=\sum_\alpha p_\alpha.
\end{equation}
We call $p_\alpha$ a \emph{log warp}.  The same expression makes sense
for a countable family whose compact supports are locally finite:
every point has a neighborhood meeting only finitely many supports.
On that neighborhood, \eqref{interface:scalar} is the scalar curvature
of a metric on its product with finitely many circles.  Adjoining a circle of
constant length one changes neither the formula nor the weight
$e^P\,\dd\vol_h$.

\begin{theorem}[Partition theorem for stretched data {\cite{OpenAIWidth}}]
\label{interface:width}
For each integer $n\ge4$, there is a finite number $D_{\mathrm{nat}}(n)>0$
with the following property.  Let $(M^n,g)$ be a connected complete smooth
Riemannian manifold without boundary, and let $D\ge D_{\mathrm{nat}}(n)$.
Choose a maximal $D$-separated set of centers $(x_i)_{i\in I}$, where
$I=\{1,\ldots,N\}$ or $I=\mathbb N$, and put
\[
 A_i=\overline B_g(x_i,D),\qquad
 B_i=M\setminus B_g(x_i,2D).
\]
Suppose every $B_i$ is nonempty.  Choose positive numbers $p_i^*$ for
$i\in\{0\}\cup I$ such that
\[
 \sum_{i\in\{0\}\cup I}p_i^*=\frac1{20}.
\]
Assume there are a complete smooth metric $h\ge g$ and a finite or
countable family of smooth real functions $(p_\alpha)$ with locally
finite compact supports such that
\begin{equation}\label{interface:width-hypothesis}
 \dist_h(A_i,B_i)>10+2\pi\sqrt{2/p_i^*}\quad(i\in I),
 \qquad R(h,(p_\alpha))\ge\frac25\quad\hbox{on }M.
\end{equation}
Then there is a continuous map $f:M\to K$ to a simplicial complex of
dimension at most $n-2$, with its usual polyhedral topology, such that
\[
 \sup_{y\in K}\diam_g f^{-1}(y)\le4D+2.
\]
The diameter bounds the entire fiber, including all its components.
\end{theorem}

We use this theorem as an imported result.  It follows from
Lemma~6.1 and Proposition~6.6 of \cite{OpenAIWidth}, followed by
Proposition~7.1 and Equation~(7.14) there.  The threshold
$D_{\mathrm{nat}}(n)$ retains the scale chosen in that paper, where
the initial scale may be enlarged.

The pointwise scalar-curvature theorem of \cite{OpenAIWidth} first
constructs $h$ and the log warps from $\Scal_g\ge1$.  Its partition
construction begins with these data already in hand.  Lemma~6.1 there
uses the separations in \eqref{interface:width-hypothesis} and the
initial scalar requirement
\[
 R(h,(p_\alpha))\ge\frac14+1.12\sum_{i\in\{0\}\cup I}p_i^*
                         =.306,
\]
which follows from $2/5$.  Its scalar hypothesis concerns
$R(h,(p_\alpha))$; a lower bound for $\Scal_g$ is not required at this
stage.

The countable version in \cite{OpenAIWidth} expressly permits local
finiteness of the initial log warps.  The measure
$e^P\,\dd\vol_h$ is smooth and positive, and every compact region
has a neighborhood on which only finitely many factors occur.
The cited construction first separates $M$ into pieces confined to
compact radial regions; label $0$ is reserved for this step.  Product
cylinders are then attached along boundary hypersurfaces, with the
metric and warp data independent of the cylinder coordinate.  The
localization and induction in Propositions~6.4 and~6.6 ensure that each
analytic piece at every later finite stage has only finitely many
relevant initial warps on its entire extent, including its product
ends.  Thus the initial data need no uniform bounds on the geometry
or the weight at infinity.  The finiteness and geometric bounds needed
on the later pieces are conclusions of the imported construction.

The reconstruction produces a continuous map $q:M\to M$ with
$d_g(x,q(x))<1$ and continuous coordinates $a_i:M\to[0,1]$ with locally
finite supports.  They satisfy
\[
 a_i=1\ \hbox{on }q^{-1}(A_i),\qquad
 a_i=0\ \hbox{on }q^{-1}(B_i),
\]
and at each point some coordinate equals $1$ while at most $n-2$
coordinates lie strictly between $0$ and $1$.  These vectors lie in a
cubical complex of dimension at most $n-2$.  Its compatible staircase
triangulations give $K$.  Local finiteness of the coordinate supports
places the map locally in a finite subcomplex, which gives continuity
for the polyhedral topology.  Equality of images preserves the full
coordinate vector.  Thus two points in one fiber have a common
coordinate equal to $1$; their $q$-images lie in the same original
outer ball, and their $g$-distance is less than $4D+2$.

To apply Theorem~\ref{interface:width} under the spectral hypothesis,
we must construct the data in \eqref{interface:width-hypothesis} at a uniform
initial stretching scale even though the positive supersolution
associated with the spectral inequality can have unbounded logarithm and
gradient.
